\documentclass[11pt]{article}

% --- dependency-light preamble (arXiv-friendly) ---
\usepackage[utf8]{inputenc}
\usepackage[T1]{fontenc}
\usepackage[margin=1in]{geometry}
\usepackage{amsmath,amssymb}
\usepackage{booktabs}
\usepackage{graphicx}
\usepackage{xcolor}
\usepackage{enumitem}
\usepackage{url}
\usepackage[hidelinks]{hyperref}

\newcommand{\code}[1]{\texttt{#1}}
\newcommand{\est}{\textsuperscript{[EST]}}
\newcommand{\tps}{tok/s}

\title{\textbf{Bit-Exact by Construction:\\
A Verification-First RTL Accelerator that Inherits the\\
GGUF $k$-Quant Checkpoint Ecosystem}}

\author{Wick-Lim\\
\small Independent Researcher\\
\small \texttt{wicklim90@gmail.com}\\
\small \url{https://github.com/Wick-Lim/WPU}}

\date{\today}

\begin{document}
\maketitle

\begin{abstract}
We present a synthesizable RTL accelerator for running a frontier open-weight
mixture-of-experts model---the 753B-parameter GLM-5.2, in its published
$\sim$467\,GB GGUF $k$-quant form---as a single-user, fully offline appliance.
The work is deliberately \emph{verification-first} and \emph{pre-silicon}: no
fabricated chip, no measured end-to-end token rate. The central claim is a
\emph{bit-exactness contract}: the dequantization of GGUF $k$-quants
(Q4\_K/Q6\_K/Q8\_0) is proven, via an independently cross-checked reference,
bitwise-equal to llama.cpp's own \code{ggml} kernels on 376{,}586{,}240 weights
from two real published GGUF files, so the silicon inherits the community's
checkpoint ecosystem by construction. This closes the ``self-referential
golden'' gap endemic to reimplementation-based accelerators. Around the
contract we contribute: (i)~speculative decoding proven in RTL---the committed
stream is a bit-exact prefix of greedy decoding (\code{spec==greedy}), enabled
by an intra-batch causal attention mode and a die-internal KV write-back, with
the amortization $A_{\mathrm{eff}}$ measured from a hardware weight-load
counter; (ii)~an adversarial verification discipline---load-bearing gates
paired with must-fail fault injections, exact per-gate test counts pinned by
the release gate, and the finding that \emph{output-insensitive} inputs (MoE
router codes absorbed by top-$k$; normalization gains) defeat end-to-end
equality checks and require direct per-beat die-input bindings; and (iii)~a
bandwidth-bound residency design whose roofline uses measured expert-union
growth and accept rates, with its stall mechanism validated on real RTL cycles
($\sim$73$\times$ under residency). The datapath places and routes on a
commodity FPGA. All projected token rates are tagged as estimates; the
remaining physical gates---vendor PHY IP and tapeout---are named, not hidden.
\end{abstract}

\section{Introduction}

Frontier large language models are increasingly available as open weights, yet
running the largest of them locally---offline, on data that must not leave the
premises---remains out of reach for most users. The obstacle is not compute but
memory bandwidth. A 753B-parameter MoE model quantized to roughly four bits
occupies $\sim$467\,GB; each decoded token reads on the order of 25\,GB of
active weights, so single-stream throughput is governed by
$\text{bandwidth} \div \text{bytes-per-token}$, and arithmetic units spend most
of their time waiting. Consumer multi-GPU rigs cannot hold 467\,GB in aggregate
VRAM, so weights spill to host memory and stream over PCIe; the effective
bandwidth---not the number of GPUs---then caps throughput at a few tokens per
second. Datacenter deployments that keep the whole model resident in
high-bandwidth memory reach hundreds of tokens per second, but at a cost and
form factor incompatible with a single-user appliance.

This paper studies the middle of that gap: an application-specific design that
holds the entire model resident in a single, wide, moderately-fast memory pool,
purpose-built to saturate the memory-bound roofline rather than to provide
general-purpose FLOPS. The target is a local box that runs the full model with
the network disconnected---an audit as literal as ``does it still work with the
Ethernet cable unplugged?''---which is precisely the property air-gapped
organizations require and which cloud, in-VPC, and trusted-execution deployments
by construction cannot provide.

\paragraph{The contract.} A design of this kind is only credible if its claims
are checkable. Machine-learning accelerators are frequently validated against a
reimplementation of the reference numerics, leaving a ``self-referential
golden'' gap: the hardware may be bit-exact to the author's own reference while
diverging from the software the community actually runs. Prior FPGA work on
GGUF-format kernels~\cite{secdallm,fbfq} accelerates $k$-quant matrix
multiplies but states no bitwise-equivalence contract with \code{ggml}; and
systems that run GGUF checkpoints on other backends via format conversion do
not in general preserve bitwise numerics. We instead make the contract the
headline: the RTL dequantization
is bit-exact to an independent reference implementation, and that reference is
proven bitwise-identical to llama.cpp's own \code{dequantize\_row\_*} kernels
on every Q4\_K and Q6\_K super-block and every Q8\_0 block of two real
published GGUF files---376{,}586{,}240 weights compared as raw \code{uint32},
no tolerance. By transitivity, the silicon dequantizes real checkpoint bytes
exactly as the software ecosystem does, so it inherits the GGUF checkpoint zoo
\emph{by construction}.
We are equally explicit about what the contract excludes: whole-runtime numeric
equality with llama.cpp (attention and accumulation orders differ by design)
and any claim about fabricated silicon.

\paragraph{Contributions.}
\begin{enumerate}[leftmargin=1.4em,itemsep=2pt]
\item \textbf{A bit-exact GGUF $k$-quant datapath, cross-validated against
llama.cpp on real checkpoint bytes} (Sections~\ref{sec:design}
and~\ref{sec:verif}). The RTL dequantization of Q4\_K/Q6\_K/Q8\_0/F16 and the
assembled GLM-5.2 forward pass (multi-head latent attention, data-dependent
sparse attention, fine-grained MoE) are bit-exact to independent goldens at the
faithful slice of Section~\ref{sec:design}; the multi-token-prediction head is
verified through the \code{spec==greedy} invariant of Section~\ref{sec:spec};
and the $k$-quant dequantization layer is additionally proven on 376.6M real
GGUF weights. To our knowledge this ecosystem-inheritance contract has not been
stated or verified by prior accelerator work.

\item \textbf{Speculative decoding proven correct in RTL} (Section~\ref{sec:spec}).
We compose a $K$-draft speculative loop into the memory-system top such that
$K{+}1$ positions are verified in \emph{one} weight load, and prove the
\code{spec==greedy} invariant: the committed stream is a bit-exact prefix of
greedy decoding under all tested accept/reject schedules. Position accuracy
requires two mechanisms absent from a shared-context batch: an
\emph{intra-batch causal attention} mode (row $j$ attends the same-pass current
keys of rows $0..j{-}1$) and a die-internal per-(layer, position) KV
write-back. The amortization $A_{\mathrm{eff}}$ is then \emph{measured from a
hardware weight-load counter}, hitting the theoretical $K{+}1$ ceiling exactly
under all-accept.

\item \textbf{An adversarial, non-vacuity verification discipline}
(Section~\ref{sec:verif}). Every load-bearing composition gate---the
speculative loop, the intra-batch oracle, each loopback family---is paired
with a fault-injection build whose failure its make target asserts; the
release gate pins the exact multiset of per-gate test counts (68 gates),
catching silently reduced testing; and default-off features are held to
structural or sequential equivalence against the pre-feature netlist where
such a gate exists (Section~\ref{sec:design}). We report a finding with, we believe,
general relevance: \emph{output-insensitive inputs defeat end-to-end equality
checks}---corrupting MoE router weight codes is absorbed by top-$k$ selection
and corrupting normalization-gain LSBs is normalized away---so proving those
bytes are consumed requires direct per-beat bindings on the die's inputs. A
``PHY-closure loopback'' applies this discipline to all five weight-input
families of the die, and a case study shows the full gate catching a real
cross-configuration elaboration regression that feature-local gates missed.

\item \textbf{A memory-bandwidth-bound residency methodology with measured
inputs} (Section~\ref{sec:analysis}). The decode roofline is parameterized by
\emph{measured} expert-union growth $U(K)$ and MTP accept rates, and the stall
mechanism it presumes is validated on real RTL cycles ($\sim$73$\times$ exposed
stall reduction under residency), while all absolute token rates remain tagged
estimates.

\item \textbf{Physical realizability evidence} (Section~\ref{sec:analysis}).
The product top synthesizes at 87.5\% LUT utilization on a commodity Kintex
UltraScale+ FPGA and a bring-up harness of it places and routes (hold met,
routed $F_{\max}$ 46.5\,MHz); the full 753B-shape hierarchy elaborates cleanly;
and a retail-parts residency prototype (1280-bit LPDDR5X, 480\,GB) is supported
by an analytical board routability study.
\end{enumerate}

\section{Background}

\paragraph{GGUF $k$-quants.} \code{llama.cpp} and its \code{ggml} tensor
library~\cite{llamacpp} are the de facto substrate for local inference, and
GGUF is their model format. The $k$-quant family stores weights in 256-element
super-blocks with a compact hierarchy of scales: Q4\_K packs 256 four-bit
weights with six-bit sub-block scales and minima into 144 bytes; Q6\_K uses 210
bytes; Q8\_0 stores 32-element blocks of eight-bit weights with an FP16 scale.
A ``dynamic'' mix such as UD-Q4\_K\_XL keeps most tensors at Q4\_K while
promoting sensitive tensors to higher precision. Matching this format
exactly---not a nearby integer quantization---is what lets an accelerator
consume the same file the community already trusts.

\paragraph{Fine-grained MoE, MLA, and MTP.} GLM-5.2~\cite{glmweights} is a
fine-grained MoE model~\cite{moe}: each token routes to a top-$k$ subset of 256
experts (plus one shared expert), so the \emph{set} of weights touched per
token is data-dependent. It uses multi-head latent attention
(MLA)~\cite{deepseekv2}, which compresses the KV cache into a low-rank latent,
a data-dependent sparse attention (DSA) index that selects top-$k$ keys per
query, and a multi-token-prediction (MTP) head~\cite{deepseekv3} usable as a
self-speculative draft in the sense of speculative
decoding~\cite{specdec,chen2023spec}.

\paragraph{Memory-bound roofline.} For single-stream decoding of a model whose
active weights exceed cache, the classic roofline model~\cite{roofline} reduces
to a bandwidth bound:
$\text{\tps} = \text{BW} \cdot A_{\mathrm{eff}} / B_{\text{tok}}$,
where BW is sustained weight bandwidth, $B_{\text{tok}}$ is bytes read per
model pass, and $A_{\mathrm{eff}}$ is the mean number of committed tokens per
pass (the speculative amortization). The design problem is to make BW large and
$B_{\text{tok}}/A_{\mathrm{eff}}$ small; the compute must merely keep up.

\section{Accelerator Design}
\label{sec:design}

\paragraph{Bit-exact datapath.} The core is a Q4\_K general matrix-multiply
(GEMM) unit that dequantizes super-blocks to FP32 and accumulates in a
documented fixed order, followed by conversion to bf16. Mixed-type columns
(Q6\_K, Q8\_0, F16) are handled by per-column type routing in the same GEMM,
fed by a weight loader that carries a per-tile type descriptor. The surrounding
operators---MLA attention with DSA key selection, the SwiGLU expert
feed-forward, the sigmoid-gated top-$k$ MoE router, RMSNorm, interleaved RoPE,
and softmax---are assembled into a full \code{glm\_model\_q4k} forward pass
(embedding $\rightarrow$ $L\times$(MLA{+}DSA{+}MoE) $\rightarrow$ final norm
$\rightarrow$ LM head $\rightarrow$ argmax). The assembled pass is bit-exact,
on logits, argmax, and hidden state, to a numpy reference that imports the same
$k$-quant dequantization used for the real-bytes cross-check of
Section~\ref{sec:verif}.

\paragraph{Slice and full shape.} Functional verification runs at a
small-but-faithful \emph{slice} that preserves every operator and ratio
(model dimension 128; six layers, three dense and three MoE; four heads;
eight experts top-$2$ plus one shared; vocabulary 256). Three further pieces of
evidence support the slice-to-full argument: (i)~the real-dimension operator
sweep, in which each operator is re-verified with its dominant dimension at the
true 753B-configuration magnitude, under the same check contract as at the
slice (the GEMM bit-exact against the \code{ggml} reference at reduction
dimension 6144; the router's top-$k$ and renormalization invariants at 256
experts top-$8$; SwiGLU, softmax, RMSNorm, and RoPE against their
high-precision goldens at intermediate 2048, length 2048, dimension 6144, and
rotary dimension 64 respectively---the GEMV reduction depth of the router
and SwiGLU stays at slice scale, that mechanism being covered by the GEMM
point); (ii)~batched-shape gates in which a
$\mathrm{PE}_M{>}1$ batched forward is proven row-wise bit-exact to standalone
single-row runs; and (iii)~a clean structural elaboration of the entire model
at the true shape (dimension 6144, 78 layers, 64 heads, 256 experts,
vocabulary 154{,}880)---a type/width check with no stimulus, which we label
\emph{elaborated}, not proven.

\paragraph{Memory system.} A single top-level module wraps the compute die with
a tier-agnostic banked read crossbar, a routed-expert cache (LRU with frequency
and prefetch), an MLA latent-KV pager with storage overflow, and the weight
loader; a two-clock wrapper adds the host/core domain-crossing FIFOs, and a
standalone, formally proven boot-load sequencer feeds the same storage path. The fabric is
\emph{byte-agnostic}---it moves addresses, slots, and IDs, never weight
bytes---so the same RTL serves streaming, full-residency, and hybrid
configurations selected by parameters. With residency enabled, expert refills
complete as tagged crossbar reads from the DRAM tier and the storage path
carries only KV spill and the one-time boot copy, an invariant checked in
simulation.

\paragraph{Result-invariant knobs.} Several parameters (a compute-resource
knob, the residency mode, the adaptive speculative depth, and every loopback
and self-KV mode of Sections~\ref{sec:spec} and~\ref{sec:verif}) are designed
to be \emph{result-invariant} or \emph{default-off}: enabling them must not
change the emitted tokens, and leaving them off must not change the netlist.
The on direction is verified for every knob by a bit-exact gate. The off
direction carries a netlist-equivalence gate where one exists: a
sequential-equivalence proof for the weight-ECC path (against an ECC-removed
reference derived from the pre-feature commit) and for boot integrity (against
the frozen pre-feature commit itself), and a structural cell-histogram
equivalence for the residency, self-KV, and CDC-protocol knobs (the first two
with a built-in liveness self-test). The loopback and adaptive-depth knobs are
default-off with behavioral gates only (bit-exact drop-in and
output-invariance); they carry no off-netlist proof. Within those stated
scopes, the knobs are optimizations, not risks.

\section{Speculative Decoding Proven in RTL}
\label{sec:spec}

Speculative decoding is the only lever that reduces the dominant roofline term:
verifying $K$ draft tokens in one model pass amortizes one weight load over up
to $K{+}1$ committed tokens. Hardware treatments of speculative
decoding~\cite{hades,specpim,specmamba} report throughput; our focus is
\emph{correctness}: a composed speculating accelerator is only trustworthy if
its committed stream provably equals what non-speculative decoding would have
produced.

\paragraph{The \code{spec==greedy} invariant.} The sequencer commits the
longest accepted \emph{prefix of the model's own argmax stream}---never a raw draft
token---so by construction the committed stream must be a bit-exact prefix of
greedy decoding, for any draft quality and any depth schedule. We verify this
end-to-end on the composed top (compute die, expert cache, KV pager, crossbars,
weight loader, and the speculative sequencer in one netlist): for
$K\in\{1,2,3\}$ under all-accept, all-reject, and mixed partial-accept
schedules, every committed token and the full logit vector of its committed row
equal a standalone non-speculative reference (31 checks).
Two independent fault-injection builds must fail and do: one commits raw
drafts instead of model argmaxes; the other writes back the KV of a
\emph{rejected} position (a phantom-KV hazard specific to partial accepts),
which diverges the very next pass's attention set.

\paragraph{Why a shared-context batch is not enough.} A naive batched verify
runs $K{+}1$ rows over a \emph{shared, already-populated} KV cache; but row
$j$, verifying the draft at position $p{+}j$, must attend all keys at positions
$0..p{+}j{-}1$, of which the last $j$ are the \emph{current-token keys of rows
$0..j{-}1$ computed in the same pass}. We therefore add, behind a default-off
parameter, an \emph{intra-batch causal attention} mode to the MLA core: each
row's current-key latent is injected into the attendable key set at causal
index $s{+}i$ (where $s$ is the cached-context length and $i$ the row index),
flowing through the same RMSNorm, up-projections, RoPE, and DSA selection as
cached keys, with the causal mask arising intrinsically from per-row key
extents. The leaf oracle proves a batched attention pass over consecutive
positions bit-exact to the serial single-row chain; the system-level batched
verify then proves full-logit bit-exactness at $K\in\{1,2,3\}$; dropping the
causal mask (a row seeing its own key) fails.

\paragraph{Die-internal KV write-back.} Position accuracy across passes
requires the die to \emph{consume KV latents it produced}. We expose the MLA
current-key latent as an output, key the pager by a flat (layer, position)
index, and close the loop: at \code{SELF\_KV=1} the die appends its computed
latent per (layer, position) and re-reads it through the pager on subsequent
steps, proven full-logit bit-exact against an independent (layer,
position)-keyed reference at one and six layers; one-off soundness checks
(documented edit-and-revert procedures in the testbench headers) confirm the
reference catches layer aliasing and append corruption; with
\code{SELF\_KV=0} the netlist matches the stub-fed base's cell histogram (the
added latent-observation ports contribute zero cells).

\paragraph{$A_{\mathrm{eff}}$ measured from a hardware counter.} The composed
top carries a \code{weight\_loads} register that increments once per
$\mathrm{PE}_M{=}K{+}1$ outer pass, and a committed-token counter. The gate
cross-checks the hardware counters against the testbench's own accounting
(a hidden per-row weight reload would fail), then reports
$A_{\mathrm{eff}} = \text{tokens}/\text{loads}$. Measured on the composed RTL:
all-accept achieves exactly the ceiling ($K{+}1$ tokens per load: 2.00, 3.00,
4.00 for $K{=}1,2,3$), all-reject the floor (1.00), and mixed schedules the
expected intermediate values. This turns the roofline's amortization term from
an assumption into a measured property of the netlist; the \emph{production}
$A_{\mathrm{eff}}$ then depends only on the model's accept rate, an externally
measured input (Section~\ref{sec:analysis}).

\section{Verification Methodology}
\label{sec:verif}

Verification is organized as a ledger. Each claim is \emph{proven} (a gated
bit-exact simulation or a solver proof), \emph{measured} (real RTL cycles or a
real place-and-route), \emph{elaborated} (structural only), or \emph{estimated}
(a roofline). Table~\ref{tab:ledger} summarizes; this section describes the
load-bearing mechanisms.

\begin{table}[t]
\centering
\small
\begin{tabular}{@{}p{0.44\linewidth}p{0.36\linewidth}l@{}}
\toprule
\textbf{Claim} & \textbf{Evidence} & \textbf{Status} \\
\midrule
$k$-quant dequant vs real GGUF & bitwise-equal to llama.cpp \code{ggml} on 376.6M weights & proven \\
Q4\_K GEMM core & bit-exact vs \code{ggml} reference (160 tests, 40 tiles) & proven \\
Assembled model forward & bit-exact vs assembled numpy golden (1{,}155 tests) & proven \\
Mixed-type path (Q6\_K/Q8\_0/F16) & bit-exact incl.\ 24-tile mixed sequence & proven \\
Speculative loop (\code{spec==greedy}) & committed stream $\equiv$ greedy prefix, 3 schedules $\times$ $K{\le}3$ & proven \\
$A_{\mathrm{eff}}$ amortization & hardware weight-load counter; all-accept $=K{+}1$/load & measured \\
Intra-batch causal attention & batch $\equiv$ serial chain, full-logit & proven \\
Die-internal KV write-back & bit-exact vs independent (layer,pos) reference & proven \\
PHY-closure loopback (5 families) & fabric round-trip bit-exact; per-beat bindings on rw/lw/gn & proven \\
Memory/control plane & BMC (7 ctrl + ECC ring) + $k$-induction (5) & proven (formal) \\
Result-invariant knobs, off-direction & cell-histogram or sequential equivalence (residency, self-KV, CDC, ECC, boot) & proven \\
Full 753B-shape hierarchy & type/width elaboration at true dimensions & elaborated \\
FPGA fit (compact config) & routed on Kintex US+, hold met, $F_{\max}$ 46.5\,MHz & measured \\
Cycle-accurate stall mechanism & per-token stall on real cycles; residency $\sim$73$\times$ & measured \\
Single-user throughput & roofline with measured $U(K)$, $r$, $A_{\mathrm{eff}}$ & estimated\est \\
End-to-end 467\,GB run; silicon rate & --- & not done \\
\bottomrule
\end{tabular}
\caption{The verification ledger. Statuses: \emph{proven} = a gated bit-exact
simulation or a solver proof; \emph{measured} = real RTL cycles or a real
place-and-route; \emph{elaborated} = structural only, no functional golden;
\emph{estimated} = a roofline projection. Every ``proven'' and ``measured''
row is reproducible from the repository; the load-bearing composition gates
are paired with fault-injection builds that must fail.}
\label{tab:ledger}
\end{table}

\paragraph{Closing the self-referential gap.} The assembled forward pass is
bit-exact to a numpy golden, but that golden is the author's own
reimplementation of \code{ggml}. To close the gap to the software the community
runs, we parse two real published GGUF files (a Qwen2.5-0.5B Q4\_K\_M and a
SmolLM2-135M Q8\_0 distribution), extract every Q4\_K and Q6\_K tensor's raw
super-blocks and every Q8\_0 tensor's blocks, and dequantize each block two
ways: with our reference, and with llama.cpp's own
\code{dequantize\_row\_q4\_K}/\code{\_q6\_K}/\code{\_q8\_0} compiled from the
built \code{ggml} library. The FP32 outputs are compared bitwise (as
\code{uint32}, not within a tolerance). All 376{,}586{,}240 weights match.
The dequantization is format-level and file-agnostic---a $k$-quant block's
bytes decode identically whatever model they came from---so bitwise identity
on these distributions' blocks transfers to any $k$-quant file, including the
467\,GB target checkpoint. By transitivity with the
RTL$\,\equiv\,$reference gates, the hardware dequantization equals the real
files' \code{ggml} dequantization. The scope is deliberate: the GEMM's FP32
accumulation follows \emph{our} documented order, not llama.cpp's runtime
order (which differs); whole-runtime equality is out of contract.

\paragraph{Injection pairing (non-vacuity).} A gate that cannot fail proves
nothing. Every load-bearing composition gate---the speculative loop (two
injections: raw-draft commit, phantom KV), the intra-batch causal oracle
(dropped mask), and each loopback family---is therefore paired with at least
one fault-injection build that must make the gate fail, and the gate's make
target asserts that failure on every run. This is mutation
analysis~\cite{mutation,mcy} specialized to one hand-chosen, semantically
meaningful mutant per claim, asserted alongside the gate itself rather than as
an offline qualification campaign; the remaining unit gates are self-checking
testbenches whose counts the manifest below pins.

\paragraph{Exact-count manifests.} Pass banners of the form ``all $N$ tests
passed'' are conventionally grepped with a wildcard $N$, which lets a
testbench silently run \emph{fewer} tests (a dropped compile flag, a reduced
loop) while CI stays green. Our release gate pins the exact multiset of test
counts for all 68 gates in a manifest; any drift---fewer tests, a vanished
configuration, an extra silent run---fails the gate. In effect the test count
itself becomes a verified contract.

\paragraph{Output-insensitive bytes and direct bindings.} While extending the
loopback proof (below) to the MoE router and normalization-gain inputs, we
found that end-to-end token equality \emph{cannot} bind them: corrupting a
router weight code by a full quantization step is absorbed by top-$k$ expert
selection (the ranking rarely changes), and corrupting a normalization gain's
LSB is normalized away downstream. A verification flow that relied only on
committed-token equality would silently accept a die that reads garbage on
those inputs. The fix is a \emph{direct per-beat binding}: hierarchical probes
assert that the die's presented input equals the expected byte for the live
request key on every beat. With the binding in place, the same corruptions that
end-to-end equality misses are caught immediately. We believe this
observation---an instance of the classic observability
problem~\cite{occom} surfacing at the integration level---is a practical
hazard for any accelerator verified primarily by end-to-end token comparison.

\paragraph{PHY-closure loopback.} Pre-silicon integration commonly stubs the
die's weight inputs from the testbench, leaving a gap: the memory fabric is
verified to \emph{count} the right transactions while the die never actually
consumes fabric-returned bytes. We close this for \emph{all five} weight-input
families of the die (attention codes, FFN expert codes, router codes, LM-head
columns, normalization gains): behind default-off parameters, each family's
bytes are routed \emph{out} of the die as tagged reads through the real banked
crossbar and back \emph{in}, clock-gating the die per beat, and the committed
stream is proven bit-exact against the stub-fed reference (thousands of
round-trips per decode; per-family non-vacuity counters). Corruption-injection
builds are asserted for the attention, expert, gain, and router
families---the latter two caught by the direct per-beat bindings that token
equality cannot see; the LM-head family is bound by the same per-beat check.
What remains external is the analog PHY itself (vendor IP); the digital
consume-path is closed.

\paragraph{Formal verification of the control plane.} The memory-system
controllers---the banked DRAM crossbar, the storage-read fabric, the boot
loader, the KV pager, the speculative-decode sequencer, the clock-throttle
controller, and the expert cache in both prefetch modes---are checked by
bounded model checking (with an
ECC-datapath variant), and five controllers are lifted to unbounded
$k$-induction, with no counterexample from reset. No formal proof touches the
numeric datapath; that plane is held by the bit-exact goldens. Residual bounded
gaps are documented rather than hidden.

\paragraph{Case study: the full gate catches what feature gates miss.} During
the intra-batch causal work, an unconditional wire computed a zero-padding
width as $(\mathrm{IDXW}{+}1)-\mathrm{DRW}$; in an unrelated top instantiating
the same MLA core with a small context and a large batch, the replication count
went negative and elaboration failed---while every gate of the \emph{feature}
passed, because the feature's own configurations never entered that corner. The
full release gate (all 68 gates, all tops) caught it. We take this as evidence
for running whole-project gates on every RTL change: in a parameterized
codebase, a shared leaf can break a sibling top at parameters the feature
author never exercises.

\section{Design-Space Analysis and Evaluation}
\label{sec:analysis}

\paragraph{Measured roofline inputs.} The bytes-per-pass denominator has two
parts that speculation amortizes differently. The routed-expert stream is
$\approx$14\,GB per pass, and under a speculative chain the \emph{union} of
experts touched across the verified positions grows by a factor $U(K)$; the
resident hot set (attention projections, dense and shared FFN, router, LM head)
is $\approx$11\,GB, read once per pass. We measure $U(K)$ directly by tracing
exact expert selections on a same-family model (GLM-4.5-Air):
$U(2){\approx}1.6$, $U(4){\approx}2.65$, $U(8){\approx}4.3$; concurrent
systems work reports compatible union growth~\cite{cascade}. We also measure
the MTP accept-rate profile $r$---first-position $r_1{\approx}0.87$ with steep
per-position decay---which places the memory-bound optimum at shallow depth
($K{=}1$--$2$) and yields a production $A_{\mathrm{eff}}{\approx}1.9$ at
$K{=}1$; the RTL-measured
mechanism of Section~\ref{sec:spec} guarantees the amortization is realized,
and the adaptive depth controller is proven output-invariant under any depth
schedule (its own gate) and designed to settle at the measured optimum.

\paragraph{Cycle-accurate validation of the mechanism.} A cycle-accurate
testbench drives the assembled system with a storage read-latency model
engaged, holding every committed token bit-exact against a standalone reference
while counting cycles. On the verified slice the compute-die baseline is
$\approx$10{,}896 cycles/token. Exposed storage stall is linear in read latency
and proportional to miss count---the memory-bound assumption, now
\emph{measured}---and the residency configuration makes stall independent of
storage latency: at a 1024-cycle latency, streaming exposes 2{,}567 stall
cycles/token whereas residency exposes 35, a $\sim$73$\times$ reduction. This
validates the residency pivot on real cycles even though the absolute product
token rate remains a projection.

\paragraph{FPGA fit.} The Q4\_K product top, in a compact configuration,
synthesizes on a Kintex UltraScale+ XCKU3P at 142{,}320 LUTs (87.5\%
utilization), $\sim$100K flip-flops, 421 DSPs, and zero BRAM; a bring-up
harness of the same top (the wide memory-side ports terminated on-chip so the
fit is not I/O-limited) then places and routes at $\sim$141.3K LUTs with hold
met and a routed maximum frequency of 46.5\,MHz, reached through
bit-exactness-preserving repipelining rounds (a campaign that raised routed
$F_{\max}$ 4.6$\times$). This establishes that the datapath is physically
realizable on commodity silicon; it is a compact slice, not the full model at
product speed, and we do not extrapolate a token rate from it.

\paragraph{Projected throughput.} Combining a resident LPDDR5X pool with the
measured $U(K)$, $r$, and the RTL-measured $A_{\mathrm{eff}}$ mechanism, the
roofline projects a residency design point of order $80$--$110$\,\tps\est\ for
the single-user box---presuming a MAC array sized to consume that bandwidth
($\sim$9.4K dequant lanes at 490\,MHz for a 1.1\,TB/s pool, $\sim$12.7K for
1.54\,TB/s, with attention and MoE holding dedicated engines that each absorb
the full bandwidth in their own sequential phase). Three cautions make this an
\emph{optimistic ceiling}: quoted bandwidths are peak, not sustained
($\sim$70--85\% is typical); measured lane scaling is sublinear
(4$\times$ lanes yielded 2.40$\times$ on the slice), so an under-provisioned
array, not memory, becomes the bound; and no board exists. For context, the
literature reports roughly 1.3--13\,\tps\ for GPU/flash offloading of models of
this scale~\cite{llminaflash,moeinfinity}, and 512\,GB-class unified-memory
workstations are community-reported around 15--25\,\tps; our figures are
projections with measured inputs, not silicon measurements.

\paragraph{Prototype design point and routability.} A physically buildable
prototype uses commodity retail parts: twenty 24\,GB LPDDR5X packages on a
1280-bit board-level bus give 480\,GB resident at $\approx$1.54\,TB/s. Because
the wide bus leaves the package, we perform a first-principles routability
study. An escape-routing model (via-in-pad HDI, the ring-cut channel-crossing
bound) estimates the $\sim$2{,}800-signal escape of the host BGA needs on the
order of four signal layers, within the six a 12-layer stack provides; a
placement check confirms the twenty packages and the host fit within a
130$\times$110\,mm board with courtyard and edge clearance. This is analysis,
not a routed board.

\section{Related Work}

\paragraph{GGUF/$k$-quant hardware.} SECDA-LLM~\cite{secdallm} integrates an
FPGA accelerator for \code{ggml} $k$-quant matrix multiplies into llama.cpp on
a PYNQ-Z1, and F-BFQ~\cite{fbfq} accelerates GGUF block quantization
(Q2\_K/Q3\_K) with dynamic format switching on a Kria KV260; both target
TinyLlama/GPT-2-class models, offload individual kernels, and state no bitwise
contract with \code{ggml}. Edge-FPGA LLM accelerators more broadly
(LlamaF~\cite{llamaf}, TeLLMe~\cite{tellme}, and full-model efforts such as
llama-fpga) use custom, ternary, or AWQ-style quantization rather than the GGUF
byte format. Our delta is not the genre but the \emph{contract}
(bitwise inheritance of the real checkpoint bytes, proven at 376.6M-weight
scale) and the \emph{scope} (a full MLA{+}DSA{+}MoE{+}MTP datapath elaborated
at 753B shape; we found no prior RTL for DSA or MTP).

\paragraph{Speculative decoding in hardware.} HADES~\cite{hades} adds
hardware-level speculative support to an LLM accelerator; SpecPIM~\cite{specpim}
co-designs speculative inference for processing-in-memory;
SpecMamba~\cite{specmamba} brings speculative decoding to FPGA for state-space
models. These report performance; none, to our knowledge, states or proves a
committed-stream-equals-greedy invariant in RTL, measures the amortization from
a hardware counter, or identifies the intra-batch causal attention and KV
write-back mechanisms required for position-accurate batched verification in a
latent-attention model.

\paragraph{MoE offload and appliances.} Software systems stream experts across
tiers (LLM-in-a-flash~\cite{llminaflash}, MoE-Infinity~\cite{moeinfinity});
architecture proposals add MoE-aware memory hierarchies
(Duplex~\cite{duplex}, Cambricon-LLM~\cite{cambricon},
Stratum~\cite{stratum}). These pursue throughput on shared or datacenter
silicon; our design targets the single-user offline residency point, and its
contribution is the verified RTL artifact and methodology rather than a new
hierarchy proposal. Concurrent measurements of expert-union growth under
multi-token verification~\cite{cascade} corroborate our $U(K)$ inputs.

\paragraph{Verification methodology.} Mutation analysis and functional
qualification (Certitude-style; MCY~\cite{mcy} in the open-source flow)
established must-fail fault builds; observability-based coverage~\cite{occom}
established that activation without propagation is invisible to output checks.
Our discipline specializes both to a per-claim, always-on CI form (one
semantically chosen mutant per gate; exact-count manifests), and documents a
concrete integration-level instance---output-insensitive weight families in an
MoE model---where end-to-end token equality is provably insufficient and
per-beat input bindings are required.

\section{Limitations and Honest Scope}

We restate the boundaries plainly. (1)~There is no fabricated silicon and no
running board; the remaining physical gates are vendor PHY IP, compiled SRAM
macros with their BIST collars, DFT insertion, and tapeout. (2)~All token rates
are roofline projections tagged \est; the FPGA result is a compact
place-and-route, not the full model at speed; measured lane scaling is
sublinear, so the projections are optimistic ceilings. (3)~Functional proofs
run at a faithful slice; the true 753B shape is structurally elaborated only
(a full-scale functional simulation is computationally infeasible), with
real-dimension operator sweeps as the bridge. (4)~The numeric contract is
\code{ggml}-exact dequantization plus a documented accumulation order;
whole-runtime equality with llama.cpp is deliberately out of contract, and the
467\,GB checkpoint has not been run end-to-end. (5)~The MTP accept rate is
measured on a same-family proxy (GLM-4.5-Air), not on GLM-5.2 itself; one
architecture parameter (\code{kv\_lora}=512) follows the DeepSeek convention
pending direct confirmation. (6)~The board study is analytical, not a routed
design. We regard stating these boundaries as part of the contribution: for a
system an auditor must sign off on, an honest ledger is the argument.

\section{Conclusion}

We have described a verification-first RTL accelerator whose defining property
is a checkable contract: the hardware dequantizes real GGUF checkpoint bytes
exactly as the community's own software does, and therefore inherits the
$k$-quant ecosystem by construction. On top of that contract we proved a
speculative decoding loop correct in RTL (committed stream $\equiv$ greedy
prefix, amortization measured from a hardware counter), formalized an
adversarial verification discipline whose gates must be able to fail---and
showed why end-to-end equality cannot see output-insensitive bytes---and
grounded a memory-bandwidth-bound residency design in measured inputs and
real-cycle stall validation. The datapath places and routes on commodity FPGA
silicon. What remains between this artifact and a device is named, not hidden:
PHY IP, macros, DFT insertion, a board, a tapeout. The full engineering record,
and every gate that backs the claims above, is open.

\section*{Reproducibility}
Every ``proven'' and ``measured'' row of Table~\ref{tab:ledger} is reproducible
from the public repository at \url{https://github.com/Wick-Lim/WPU} via the
corresponding build targets: the bit-exact simulation gates and their paired
injection builds, the GGUF cross-check tool, the exact-count release manifest
(\code{make release-gate-strict}), the speculative and loopback gates, the
formal model-checking targets, and the cycle-accurate performance harness.

\section*{Acknowledgments}
The GGUF cross-check builds on the open-source \code{ggml}/llama.cpp project
and the published GLM open weights; this work is an independent inference
accelerator for those models.

\begin{thebibliography}{99}

\bibitem{roofline}
S.~Williams, A.~Waterman, and D.~Patterson.
\newblock Roofline: an insightful visual performance model for multicore
architectures.
\newblock \emph{Communications of the ACM}, 52(4):65--76, 2009.

\bibitem{moe}
N.~Shazeer, A.~Mirhoseini, K.~Maziarz, A.~Davis, Q.~Le, G.~Hinton, and J.~Dean.
\newblock Outrageously large neural networks: the sparsely-gated
mixture-of-experts layer.
\newblock In \emph{ICLR}, 2017. arXiv:1701.06538.

\bibitem{specdec}
Y.~Leviathan, M.~Kalman, and Y.~Matias.
\newblock Fast inference from transformers via speculative decoding.
\newblock In \emph{ICML}, 2023. arXiv:2211.17192.

\bibitem{chen2023spec}
C.~Chen, S.~Borgeaud, G.~Irving, J.-B.~Lespiau, L.~Sifre, and J.~Jumper.
\newblock Accelerating large language model decoding with speculative sampling.
\newblock arXiv:2302.01318, 2023.

\bibitem{deepseekv2}
DeepSeek-AI.
\newblock DeepSeek-V2: a strong, economical, and efficient mixture-of-experts
language model.
\newblock arXiv:2405.04434, 2024.

\bibitem{deepseekv3}
DeepSeek-AI.
\newblock DeepSeek-V3 technical report.
\newblock arXiv:2412.19437, 2024.

\bibitem{llamacpp}
G.~Gerganov and the llama.cpp/ggml contributors.
\newblock llama.cpp and the ggml tensor library.
\newblock Software, \url{https://github.com/ggml-org/llama.cpp}.

\bibitem{glmweights}
Zhipu AI / THUDM.
\newblock GLM-5.2 open weights (GGUF $k$-quant distribution).
\newblock Model weights, \url{https://huggingface.co/unsloth/GLM-5.2-GGUF}.

\bibitem{secdallm}
J.~Haris, R.~Saha, W.~Hu, and J.~Cano.
\newblock Designing efficient LLM accelerators for edge devices (SECDA-LLM).
\newblock arXiv:2408.00462, 2024.

\bibitem{fbfq}
J.~Haris and J.~Cano.
\newblock F-BFQ: flexible block floating-point quantization accelerator for
LLMs.
\newblock arXiv:2510.13401, 2025.

\bibitem{llamaf}
H.~Xu, Y.~Li, and S.~Ji.
\newblock LlamaF: an efficient Llama2 architecture accelerator on embedded
FPGAs.
\newblock arXiv:2409.11424, 2024.

\bibitem{tellme}
Y.~Qiao, Z.~Chen, Y.~Zhang, Y.~Wang, and S.~Huang.
\newblock TeLLMe v2: an efficient end-to-end ternary LLM prefill and decode
accelerator with table-lookup matmul on edge FPGAs.
\newblock arXiv:2510.15926, 2025.

\bibitem{hades}
Z.~Yang, Y.~Jin, and X.~Xu.
\newblock HADES: hardware accelerated decoding for efficient speculation in
large language models.
\newblock arXiv:2412.19925, 2024.

\bibitem{specpim}
C.~Li et al.
\newblock SpecPIM: accelerating speculative inference on PIM-enabled systems
via architecture-dataflow co-exploration.
\newblock In \emph{ASPLOS}, 2024.

\bibitem{specmamba}
L.~Zhong et al.
\newblock SpecMamba: accelerating Mamba inference on FPGA with speculative
decoding.
\newblock arXiv:2509.19873, 2025.

\bibitem{llminaflash}
K.~Alizadeh et al.
\newblock LLM in a flash: efficient large language model inference with limited
memory.
\newblock arXiv:2312.11514, 2023.

\bibitem{moeinfinity}
L.~Xue et al.
\newblock MoE-Infinity: efficient MoE inference on personal machines with
sparsity-aware expert cache.
\newblock arXiv:2401.14361, 2024.

\bibitem{duplex}
S.~Yun et al.
\newblock Duplex: a device for large language models with mixture of experts,
grouped query attention, and continuous batching.
\newblock In \emph{MICRO}, 2024.

\bibitem{cambricon}
Z.~Yu et al.
\newblock Cambricon-LLM: a chiplet-based hybrid architecture for on-device
inference of 70B LLM.
\newblock In \emph{MICRO}, 2024. arXiv:2409.15654.

\bibitem{stratum}
Y.~Pan et al.
\newblock Stratum: system-hardware co-design with tiered monolithic
3D-stackable DRAM for efficient MoE serving.
\newblock In \emph{MICRO}, 2025. arXiv:2510.05245.

\bibitem{cascade}
A.~Saxena, P.-A.~Tsai, H.~Taneja, A.~Jaleel, and M.~Qureshi.
\newblock Utility-driven speculative decoding for mixture-of-experts.
\newblock arXiv:2506.20675, 2025.

\bibitem{mutation}
Y.~Jia and M.~Harman.
\newblock An analysis and survey of the development of mutation testing.
\newblock \emph{IEEE Transactions on Software Engineering}, 37(5):649--678,
2011.

\bibitem{mcy}
YosysHQ.
\newblock MCY: mutation cover with Yosys.
\newblock Software, \url{https://github.com/YosysHQ/mcy}.

\bibitem{occom}
F.~Fallah, S.~Devadas, and K.~Keutzer.
\newblock OCCOM: efficient computation of observability-based code coverage
metrics for functional verification.
\newblock In \emph{DAC}, 1998.

\end{thebibliography}

\end{document}
